% !PRINTED 397
\weberpage{397}
%<<CONT
be, because 2), as a real complex, cannot contain a real double tangent
paired with an imaginary one.

Let then
\[
2)\quad \xi\xi',\ x_1x_2,\ y_1y_2,
\]
and from this we derive the two complexes:
\[
\begin{aligned}
3)&\quad x_1\xi,\ x_2\xi'\\
4)&\quad x_1\xi',\ x_2\xi,
\end{aligned}
\]
which are syzygetic with one another, and hence have no element in
common besides the four indicated. But they are at the same time
conjugate, and therefore no real double tangent besides $x_1, x_2$ can
occur in 3). Now, however, 1) and 3) are azygetic, because
$\xi$ does not occur in 1); hence one element from each pair of 1)
must occur in 3). Thus 3) contains real elements besides $x_1, x_2$,
whereby the contradiction results.

We conclude from this that the two complexes which are determined
by the pairs $x_1x_2$, $y_1y_2$ and which form with 1) a syzygetic
triple of complexes contain \gesperrt{only real elements}.

If we let one of these complexes, say $x_1x_2$, take the place
of 1) and then repeat our argument, it follows
that all the double tangents whatsoever, and so also those of the
last pair of 1), must be real.

5. \gesperrt{If more than eight real double tangents
are present, then there are sixteen which form
four real pairs each in a syzygetic triple of
complexes.}

If we assume more than eight real double tangents, then by
the theorems 3., 4. we can assemble the following syzygetic triple
of complexes:
\[
\begin{aligned}
1)&\quad x_1y_1,\ x_2y_2,\ x_3y_3,\ x_4y_4\\
2)&\quad x_1x_2,\ y_1y_2,\ z_1z_2\\
3)&\quad x_1y_2,\ x_2y_1,
\end{aligned}
\]
in which the $x_i, y_i, z_i$ are real.

With this we combine the complex
\[
4)\quad x_1z_1,\ x_2z_2,
\]
which is azygetic with 1), and can therefore contain one and only
one element from each of the pairs $x_3y_3, x_4y_4$, say $x_3$ and $x_4$.
But since 4) is a real complex, there must be two further
%CONT>>
